\chapter{Practical Demonstrator}
\label{chap:demonstrator}

Chapter~\ref{chap:evaluation} evaluates the model on a held-out slice drawn from
the same seed distribution it was trained on. This chapter asks a different and
harder question: what does the model do when applied to repositories it was never
sampled to see? Answering it serves the demonstration phase of the Design Science
Research process \cite{hevner2004design,peffers2007design}, but its scientific
value here is as an error analysis. An out-of-distribution application exposes
failure modes that a held-out slice cannot, because the held-out slice inherits
the seed's composition and therefore its blind spots.

The input is the dependency set of a real project, scored through the deployed
artifact. Because the staging gate guarantees that the deployed model is the
evaluated model (Section~\ref{sec:impl-deployment}), the outputs below are
evidence about the same \texttt{xgboost\sdash full\sdash history} artifact (feature set
\texttt{feature\sdash set\sdash v4\sdash full\sdash history}) whose metrics Chapter~\ref{chap:evaluation}
reports, and not about a separately configured system. All figures are real
outputs and reflect the public repository state on \texttt{2026-07-09}, the date
of the analysis run.

\section{Example Repository Set}
\label{sec:demo-mapping}

Consistent with the unit of analysis fixed in Section~\ref{sec:impl-architecture},
each repository is scored independently and the dependency inventory is taken as an
external input produced by software-composition-analysis tooling
\cite{ortAnalyzer}.

For this demonstration the input is the set of dependency repositories of
\texttt{ArchipelagoMW\sslash Archipelago}, a large and actively developed open-source
multi-world game randomizer. Its declared Python dependencies correspond to 30
source repositories, ranging from widely used infrastructure libraries (for
example \texttt{flask}, \texttt{jinja2}, \texttt{werkzeug}) to smaller, more
specialized packages. This project was chosen because that mixture is expected to
produce a non-trivial triage view rather than a uniformly low-risk one.

Each repository is scored through the deployed analysis workflow, which enriches
live repository metadata and applies the staged model ensemble. The two prediction
regimes of Section~\ref{sec:method-prediction-regimes} appear directly in this
run: one repository (\texttt{flask}) is present in the staged repository feature
cache and is therefore scored in the richer full-history regime, while the
remaining 29 are absent from the foundation cache and fall back to the cold-start
regime. This distribution is itself an honest illustration of the cold-start
coverage limitation highlighted in Chapter~\ref{chap:conclusion}: an arbitrary
project's dependency repositories are mostly outside the precomputed foundation
seed. The aggregate held-out gap between the two regimes is small, but
Section~\ref{sec:results-history-ablation} shows that this understates the cost of
scoring without reconstructed history: on repositories still active at the
observation time --- which is what most of the thirty are --- the full-history
model outranks the cold-start one by \(0.030\) ROC-AUC. Twenty-nine of these
thirty scores are therefore drawn from the weaker of the two regimes, on precisely
the population where the difference is largest. A second caveat compounds it: the
cold-start metrics are measured on seed-distribution repositories while the model
is applied here to arbitrary submissions.

\section{Risk Ranking Output and Interpretation Guidance}
\label{sec:demo-ranking}

The scoring service produces, for each repository, the inactivity-risk score
(0--100), the discrete action level (\texttt{monitor}, \texttt{review}, or
\texttt{replace\_candidate}), a confidence value derived from signal completeness
and prediction decisiveness, and the prediction regime. Of the 30 scored
repositories, the risk buckets are two critical, three high, eight medium, and 17
low; by action level, two repositories are raised to
\texttt{replace\_candidate}, four to \texttt{review}, and the remaining 24 to
\texttt{monitor}. The ranking therefore concentrates analyst attention on six of
thirty repositories rather than diluting it across the whole set.
Table~\ref{tab:demo-ranking} shows the flagged repositories together with the
full-history item for contrast; the complete 30-repository ranking is given in
Appendix~\ref{appendix:demo-full}. Every one of the 30 scores is distinct, which
reflects the interpolated calibration mapping of
Section~\ref{sec:impl-models} rather than the flat per-bin lookup it replaced.

\begin{table}[htbp]
  \centering
  \small
  \caption{Flagged dependency repositories of \texttt{ArchipelagoMW/Archipelago}
  (the two \texttt{replace\_candidate} items and the four \texttt{review} items),
  with the single full-history repository shown for contrast. Risk is the 0--100
  inactivity-risk score; ``Conf.'' is the per-prediction confidence; ``Regime'' is
  full-history (FH) or cold-start (CS). The complete 30-repository ranking is in
  Appendix~\ref{appendix:demo-full}.}
  \label{tab:demo-ranking}
  \setlength{\tabcolsep}{4pt}
  \begin{tabular}{llrlrl}
    \toprule
    Package & Source repository & Risk & Action & Conf. & Regime \\
    \midrule
    bsdiff4      & \texttt{ilanschnell/bsdiff4}       & 92.1 & \texttt{replace\_candidate} & 0.69 & CS \\
    setproctitle & \texttt{dvarrazzo/py-setproctitle} & 80.4 & \texttt{replace\_candidate} & 0.61 & CS \\
    \midrule
    jinja2       & \texttt{pallets/jinja}             & 77.1 & review  & 0.77 & CS \\
    cymem        & \texttt{explosion/cymem}           & 74.9 & review  & 0.75 & CS \\
    markupsafe   & \texttt{pallets/markupsafe}        & 74.2 & review  & 0.75 & CS \\
    pymem        & \texttt{srounet/Pymem}             & 54.3 & review  & 0.47 & CS \\
    \midrule
    flask        & \texttt{pallets/flask}             &  9.5 & monitor & 0.91 & FH \\
    \bottomrule
  \end{tabular}
\end{table}

\paragraph{The highest-ranked repository.} The item the ranking elevates most
strongly is \texttt{bsdiff4} (\texttt{ilanschnell\sslash bsdiff4}), a small binary-diff
library, with a risk score of \(92.1\). Its explanation factors, taken directly
from the scored output, are consistent with the intended semantics of the model:
the two cold-start ensemble members predict a \(97.0\%\) and \(87.2\%\) inactivity
risk respectively, the model-agreement factor reports the spread between them, and
the input-completeness factor records that \(75\%\) of expected signals were
available before imputation. The score reflects a genuinely quiet maintenance
profile --- a low-popularity, single-maintainer repository with no detected release
cadence --- which is exactly the kind of dependency a supply-chain reviewer would
want surfaced. The confidence of \(0.69\) is moderate rather than high, and caveats
record both that the score was produced by the cold-start model, because no
full-history cache row was available, and that two release-related signals had to
be imputed. The flag is therefore presented as a prompt for manual review rather
than as a definitive judgement.

\paragraph{Two instructive false positives.} The ranking also raises
\texttt{jinja2} (\(77.1\)) and \texttt{markupsafe} (\(74.2\)) to \texttt{review}.
Both are Pallets projects: widely used, professionally maintained, and in no
meaningful sense at risk of abandonment. Neither elevation can be blamed on missing
data --- both were scored with \(100\%\) signal completeness --- so these are
genuine model errors rather than artifacts of imputation, and they are reported
here rather than omitted. The explanation is the failure mode that
Section~\ref{sec:background-metric-pitfalls} anticipates: a mature, feature-complete
library that has reached a stable API exhibits precisely the observable profile of
a project winding down --- infrequent commits, long gaps between releases, a small
active contributor set --- because low activity and low \emph{need} for activity are
indistinguishable in public metadata. The model was trained on a label that defines
inactivity as the absence of observable maintenance signals
(Section~\ref{sec:method-definitions}), so it is behaving exactly as specified;
what it cannot express is that for \texttt{jinja2} the absence of churn is a sign of
maturity. This is the strongest practical argument in this thesis for treating the
score as a prioritization signal rather than a verdict, and for the human review
step that the action levels are designed to trigger.

\paragraph{A low-risk repository for contrast.} At the opposite end,
\texttt{flask} (\texttt{pallets/flask}) receives a risk score of \(9.5\) with a
confidence of \(0.91\). It is the one repository scored in the full-history regime,
with complete signal coverage and no imputation, and its two ensemble members agree
closely (\(5.4\%\) and \(13.5\%\) inactivity risk). That \texttt{flask} is ranked
low while its sibling projects \texttt{jinja2} and \texttt{markupsafe} are flagged
is instructive in itself: the three are maintained by the same organization under
the same release practices, and the difference in their scores tracks the
availability of reconstructed history rather than any difference in maintenance.
This contrast between a moderate-confidence cold-start flag and a
high-confidence full-history clearance illustrates why the confidence value, not
only the risk score, is part of the triage output --- and it quantifies the
cold-start coverage gap more sharply than the held-out metrics do.

Interpretation guidance follows the rest of the thesis. The score is a
prioritization signal, not a trust or security verdict: a low score does not
certify a repository as secure, and an elevated score indicates thin public
maintenance signals rather than a known defect. Low confidence or a large number of
imputed signals should reduce the decisiveness of any action, and cold-start scores
in particular are conditional on the reduced feature set and on the sampled
prevalence to which the model is calibrated
(Section~\ref{sec:method-limitations}), so they warrant recalibration before being
read as absolute deployment probabilities. Used this way, the ranked view converts
a flat list of 30 repositories into six candidates for review, of which at least
two --- on the evidence above --- a human reviewer would quickly clear. That is a
useful reduction of analyst effort, and an honest statement of the artifact's
precision at this operating point.
